\documentclass[10pt,a4paper]{amsart}
\usepackage[margin=19mm]{geometry}
\usepackage{amsmath,amssymb,mathtools}
\usepackage[scr=rsfs,cal=euler]{mathalfa}
\usepackage{xfrac}
\usepackage[unicode,hidelinks,bookmarks=false]{hyperref}
\newcommand{\ie}{i.e.\ }
\newcommand{\cf}{cf.\ }
\newcommand{\resp}{respectively\ }
\newcommand{\Subsection}{\subsection}
\providecommand{\nobreakdash}{\nobreak\hbox{-}}
\setlength{\emergencystretch}{2em}
\multlinegap0pt
\usepackage{bm}
\usepackage{bbm}
\usepackage{dsfont}
\usepackage{xspace}
\usepackage{cancel}
\usepackage{mathtools}
\usepackage{cleveref}
\usepackage{amsthm}
\makeatletter
\@ifundefined{theorem}{\newtheorem{theorem}{Theorem}}{}
\@ifundefined{assumption}{\newtheorem{assumption}[theorem]{Assumption}}{}
\@ifundefined{definition}{\newtheorem{definition}[theorem]{Definition}}{}
\@ifundefined{lemma}{\newtheorem{lemma}[theorem]{Lemma}}{}
\@ifundefined{notation}{\newtheorem{notation}[theorem]{Notation}}{}
\makeatother
\DeclareMathOperator{\Spec}{Spec}
\DeclareMathOperator{\Trop}{trop}
\DeclareMathOperator{\trop}{trop}
\newcommand{\RR}{\mathbb{R}}
\newcommand{\an}{{\mathrm{an}}}
\title[Generic root counts and flatness in tropical geometry]{Generic root counts and flatness in tropical geometry}
\author{Paul Alexander Helminck, Yue Ren}
\date{Source: arXiv:2206.07838v3 (CC BY 4.0)}
\begin{document}
\maketitle

\noindent\emph{Source and licence.} Generic root counts and flatness in tropical geometry. Version arXiv:2206.07838v3 (CC BY 4.0), distributed under the Creative Commons Attribution 4.0 International licence, as recorded in the arXiv licence metadata for this exact version and shown on the abstract page https://arxiv.org/abs/2206.07838v3. Full rights notice: licenses/arxiv-2206.07838-CC-BY-4.0.txt.\par\smallskip

\noindent\textit{Editorial scope.} Counted unit: the Theorem whose source label is \texttt{thm:TropicalGenericFlatness2} in the article (source0.tex lines 1265--1268, source label \texttt{thm:TropicalGenericFlatness2}), delivered together with the article's own complete original proof of it (source0.tex lines 1269--1284, one \texttt{proof} environment of 16 lines). No mathematics is written, repaired or re-derived: statement and proof are the article's own text line for line. Required context retained uncounted: con:mainConvention (lines 478--486); def:TropicalBases (lines 919--932); def:TropicallyFlat (lines 1002--1005); lem:GenericValuations3 (lines 1132--1134); lem:FindingRationalPoints (lines 1147--1149); lem:ExtendValuationPoints (lines 1216--1226); assumption:witnessSet (lines 1240--1245); lem:fiberwiseTropicalizationLocallyConstant (lines 1247--1252). Every definition, display and label of those lines is retained unchanged. Omitted from this package and not redistributed: 1--477, 487--918, 933--1001, 1006--1131, 1135--1146, 1150--1215, 1227--1239, 1246--1246, 1253--1264, 1285--2259. Nothing inside a retained excerpt is omitted or altered: every retained body is delivered exactly as the article's lines are written, so no omission receipt of any kind accompanies this package. Citations: the retained excerpts carry no citation command at all, so no bibliography entry of the article is transcribed and no reference is redistributed. Reference adaptation: none was needed and none was made; every reference inside the delivered unit resolves inside it, so no unresolved reference or citation is produced. Printed slips kept literal: no printed word, symbol, hypothesis or number of the counted statement or of its proof was altered. Layout and class: the article's own class is \texttt{amsart} with packages including biblatex, hyperref, inputenc, amsmath, amsthm, amssymb, amsfonts, enumerate, mathtools, tikz-cd, hhline, stmaryrd, enumitem, centernot; the wrapper is the lane's pinned \texttt{amsart} editorial wrapper at 10pt on A4, which restates the theorem-like environments the retained text needs and adds the article's own macro definitions that the retained text needs (\texttt{\textbackslash RR}, \texttt{\textbackslash Spec}, \texttt{\textbackslash Trop}, \texttt{\textbackslash an}, \texttt{\textbackslash trop}), with extra wrapper packages (bm, bbm, dsfont, xspace, cancel, mathtools, cleveref). The delivered numbering is the unit's own. Assets: no retained span contains a figure environment, a graphics-inclusion command, a PDF-page inclusion, a table or a TikZ picture; no image file is copied into this package, the package declares no asset dependency and no image file is redistributed with it. Licence: version arXiv:2206.07838v3 of this article is distributed under the Creative Commons Attribution 4.0 International licence, as recorded in the arXiv licence metadata for this exact version and shown on the abstract page https://arxiv.org/abs/2206.07838v3. A full rights notice is delivered at licenses/arxiv-2206.07838-CC-BY-4.0.txt. Characters: every retained excerpt is pure ASCII, so the delivered engine, XeLaTeX with its default Latin Modern text font, reports no missing character.
\begin{notation}\label{con:mainConvention}
  For the remainder of the paper, let $K$ be an algebraically closed field with a non-archimedean absolute value $|\cdot|_K\colon K\rightarrow\RR_{\geq 0}$.

  Let $A$ be a $K$-algebra of finite type and let $Y=\Spec(A)$ be its associated scheme.  We will refer to $A$ as the \emph{parameter ring} and $Y$ as the \emph{parameter space}.  We will assume $Y$ to be integral and use $\eta$ to denote its unique generic point.  Moreover, abbreviating $A[x^\pm]\coloneqq A[x_1^\pm,\dots,x_n^\pm]$, let $T\coloneqq \Spec(A[x^\pm])$ be an $n$-dimensional torus over $Y$, and denote the projection by $p\colon T\rightarrow Y$. If $P\in Y$, then we write $k(P)$ for the residue field of $P$. There is a natural ring homomorphism $A\to k(P)$ and for $f\in{A}$, we write $f(P)\in k(P)$ for the image of $f$ under this homomorphism.

  Let $B=A[x^\pm]/I$ for some ideal $I\subseteq A[x^\pm]$. We identify $X\coloneqq \Spec(B)$ with a closed subspace of $T=\Spec(A[x^{\pm}])$ through the closed immersion induced by the natural ring homomorphism $A[x^{\pm}]\to A[x^{\pm}]/I$. We will not distinguish between $X$ and its image in $T$.  %
  By composing the inclusion $X\to T$ with $p$, we obtain a morphism $p_{X}\colon X\to Y$, and we will often abbreviate $p=p_X$ if the context is clear.
  For a point $P\in Y$, we denote $A[x^\pm]_P\coloneqq A[x^\pm]\otimes_K k(P)$ and $T_P\coloneqq \Spec(A[x^\pm]_P)$ as well as $B_P\coloneqq B\otimes_K k(P)$ and $X_P\coloneqq \Spec(B_P)$. We refer to $X_P$ as the \emph{specialization} of $X$ at $P$.
\end{notation}

\begin{definition}\label{def:TropicalBases}
  Let $P\in Y^\an$, and let $V\subseteq Y$ be a Zariski-open set with $\pi(P)\in V$. Let $A_V$ be the induced coordinate ring, and $I_V$ be the induced ideal.  Let $f_1,\dots,f_k\in I_V\subseteq A_V[x^\pm]$ be a set of generators, say $f_i=\sum c_{i,\alpha}x^{\alpha}$ for some $c_{i,\alpha}\in A_V$.

  We say $f_1,\dots,f_k$ are a \emph{$K$-rational local tropical basis} of $X$ around $P$ if there is an open neighbourhood $U\subseteq Y^\an$ with $P\in U$ and $\pi(U)\subseteq V$ such that $f_{1,Q},\dots,f_{k,Q}\in I_Q$ is a tropical basis for all $Q\in U$, i.e., $\Trop(X_{Q})=\bigcap_{i=1}^{k} \Trop(V(f_{i})_{Q})$. Moreover, we say $f_1,\dots,f_k$ are \emph{non-degenerate} around $P$, if $c_{i,\alpha}(Q)\neq 0$ for all $Q\in U$ unless $c_{i,\alpha}(P)=0$.

  A \emph{(non-degenerate) local tropical basis} around $P$ is a (non-degenerate) $L$-rational local tropical basis of $X_L$ at $P_L$ for some valued field extension $K\to{L}$ and a point $P_{L}$ lying over $P$.
 %
 %
 %
 %
 %
 %
 %
\end{definition}

\begin{definition}
  \label{def:TropicallyFlat}
  We say $X$ is \emph{tropically flat} over $P\in{Y^{\mathrm{an}}}$ if it admits a non-degenerate local tropical basis at $P$ as defined in \cref{def:TropicalBases}. The \emph{tropically flat locus} is the set of all $P\in{Y^{\mathrm{an}}}$ over which $X$ is tropically flat.
\end{definition}

\begin{lemma}\label{lem:GenericValuations3}
  Any basic open set $\mathbb{B}(r_{1},r_{2},h)\coloneqq \{Q\in Y^\an\mid r_1<|h(P)|<r_2\}$, where $h\in {A\backslash K}$ and $0\leq r_{1}<r_{2}$, contains a generic valuation $P\in Y^{\an}$.
\end{lemma}

\begin{lemma}\label{lem:FindingRationalPoints}
  Let $P$ be a generic valuation and let $\mathcal{C}\subseteq A$ be a finite non-empty subset. Suppose that $K$ is non-trivially valued and algebraically closed. Then there is a rational point $Q\in \mathbb{B}_{P,\mathcal{C}}$ and an open neighborhood $U\subseteq \mathbb{B}_{P,\mathcal{C}}$ around $Q$.
\end{lemma}

\begin{lemma}\label{lem:ExtendValuationPoints}
  Let $P\in Y^{\an}$ be a generic valuation. For any finite number of tropical points $w_{1},\dots, w_{k}\in\trop(X_{P})\cap \mathrm{val}(k(P))^n$, there exists a finite extension of valued fields $k(P)\to L$ and points $z_1,\dots,z_k\in X_{P}(L)\subseteq L^n$ such that
  \begin{math}
    \mathrm{val}(z_{j})=w_{j},
  \end{math}
  with $\mathrm{val}(\cdot)$ denoting coordinatewise valuation.
  %
  %
  %
  %
\end{lemma}

\begin{assumption}
  \label{assumption:witnessSet}
  For the remainder of the section, fix a witness set $Z=\{z_1,\dots,z_m\}$ $\subseteq X_P(L)$.  Let $A'$ be the integral closure of $A$ in $L$. This defines a finite normalization morphism $\mathrm{norm}\colon Y'\coloneqq \Spec(A')\to Y=\Spec(A)$, and the extended valuation from \cref{lem:ExtendValuationPoints} directly gives a point $P'\in Y'^{\an}$ mapping to $P\in Y^{\an}$.  Moreover, we can find an open neighborhood $V'\subseteq Y'$ such that $z_{i,j}\in A'_{V'}$. %

  Note that, by Grothendieck's generic flatness theorem, there is open subset $V\subseteq Y$ such that $\mathrm{norm}^{-1}(V)\to V$ is flat and thus open.  By restricting to open subsets of $Y'$ and $Y$, we may therefore assume that $z_{i,j}\in A'$.
\end{assumption}

\begin{lemma}\label{lem:fiberwiseTropicalizationLocallyConstant}
  Let $K$ be non-trivially valued, and let $I\subseteq A'[x_1,\dots,x_n]$.  Let $P\in{Y'^{\mathrm{an}}}$ be a generic valuation. Then there is a finite subset $\mathcal{C}\subseteq A'$ such that for all $Q\in \mathbb{B}_{\mathcal{C}}$
  \begin{equation*}
    \mathrm{trop}(V(I)_{P})=\mathrm{trop}(V(I)_{Q}).
  \end{equation*}
\end{lemma}

\begin{theorem}\label{thm:TropicalGenericFlatness2}
  Let $X=V(I)$ and $Y$ be as in \cref{con:mainConvention}.
  Then the tropically flat locus of $X$ contains a dense open subset of $Y^{\mathrm{an}}$. If $K$ is non-trivially valued, then this locus moreover contains a dense open subset of $Y(K)\subseteq Y^\mathrm{an}$.
\end{theorem}

\begin{proof}
  Since we allow valued field extensions in \cref{def:TropicallyFlat} of tropically flatness, we can assume that $K$ is non-trivially valued.

  We first show the statement for $Y'=\Spec(A')$ and $X'=X\times_{Y} Y'$ from \cref{assumption:witnessSet}. Consider a basic open set $\mathbb{B}(r_{1},r_{2},h)\subseteq Y'^{\mathrm{an}}$ for some $r_1,r_2\in\RR_{\geq 0}$ and $h\in{A'}$.
  Using \cref{lem:GenericValuations3}, we can find a generic valuation $P\in \mathbb{B}(h,r_{1},r_{2})$.  Let $f_1,\dots,f_m\in A'[x_1,\dots,x_n]$ such that $f_{1,P},\dots,f_{m,P}$ is a tropical basis of $I_P$, and let $\mathcal C_f\subseteq A'$ be the set of coefficients of $f_1,\dots,f_m$.
  By \cref{lem:fiberwiseTropicalizationLocallyConstant}, there is a $\mathcal C_{\trop}\subseteq A'$ such that $\trop(X'_P)=\trop(X'_Q)$ for all $Q\in\mathbb{B}_{\mathcal{C}_{\trop}}$.
  Set $\mathcal C\coloneqq \{h\}\cup\mathcal C_f\cup\mathcal C_{\trop}$.
  By \cref{lem:FindingRationalPoints}, there is a rational point $Q_{0}\in \mathbb{B}_{\mathcal{C}}$ and an open neighborhood $Q_{0}\in U\subseteq \mathbb{B}_{\mathcal{C}}$.
  We find that $X$ is tropically flat around any point $Q\in{U}$. We conclude that the tropically flat locus of $X'$ contains an open and dense subset of $Y'$ as well as an open and dense subset of $Y'(K)$.

  We now treat the general case.  Consider a basic open set $\mathbb{B}\subseteq Y^\an$ containing a generic valuation $P$.  Let $\mathbb{B}'$ be the preimage of $\mathbb{B}$ under the open normalization map $\phi\colon Y'\to Y$ from \cref{assumption:witnessSet}.  Note that the point $P'$ from \cref{assumption:witnessSet} is in~$\mathbb{B}'$. Using what we proved above, we find an open neighborhood $U'\subseteq \mathbb{B}'$ of $P'$ such that $X'$ is tropically flat over $U'$.  Set $U=\phi(U')$, which is open as $\phi$ is open.  For all $Q'\in U'$ and $Q=\phi(Q')$, we have
  \begin{equation*}
\trop(X'_{Q'})=\trop(X_{Q})=\trop(X'_{P'})=\trop(X_{P}).
  \end{equation*}
  We thus find that $X$ is tropically flat over $U\subseteq \mathbb{B}$. In particular, the tropically flat locus contains an open and dense subset.
\end{proof}

\end{document}
